% Symbolic threshold analysis for the synthetic BM-JEPA experiments.
% This derivation matches data/synthetic.py.

\paragraph{Generator geometry.}
Let $r=r_\star$ and let $\{e_i\}_{i=0}^{r-1}$ denote the latent coordinate
basis.  The revised generator separates target heterogeneity from target
overlap.  First, it constructs a dictionary of possible target directions
\[
    v_0=e_0,\qquad
    v_i(\alpha)
    =
    \sqrt{1-h_i(\alpha)^2}\,e_0+h_i(\alpha)e_i,
    \qquad i=1,\ldots,r-1,
\]
where
\[
    h_i(\alpha)=\alpha^{p_i},
    \qquad
    p_i\in[0.75,1.65].
\]
The direction-dependent exponents avoid making every non-shared singular value a
trivial linear function of the sweep parameter.  Nevertheless, $\alpha$ still
has the intended interpretation: small $\alpha$ makes all atoms nearly
collinear with $e_0$, while $\alpha=1$ recovers the coordinate basis.

\paragraph{Target-bank coverage.}
Let $q_i$ be the number of target rows assigned to atom $i$ by the base balanced
assignment.  The overlap parameter does not change the dictionary geometry.
Instead, it changes the coverage distribution over atoms.  Define
\[
    \pi_i(\omega)
    =
    \frac{(q_i/K)\exp[-\kappa(\omega)t_i]}
         {\sum_{\ell=0}^{r-1}(q_\ell/K)\exp[-\kappa(\omega)t_\ell]},
    \qquad
    t_i=\frac{i}{r-1},
\]
with
\[
    \kappa(\omega)=\frac{4\omega^2}{1-\omega}.
\]
The per-row scale for atom $i$ is chosen so that the aggregate energy assigned
to that atom is proportional to $\pi_i(\omega)$.  Thus $\omega=0$ preserves the
balanced target bank, while larger $\omega$ concentrates coverage on the shared
low-index atoms.  This is a redundancy/coverage control rather than another
angular interpolation.

\paragraph{Population spectrum.}
Ignoring target noise, the stacked target operator is $A\in\mathbb{R}^{K\times r}$.
Since the context operator has orthonormal columns, the nonzero singular values
of the population cross-covariance $\Sigma_{YZ}=AC^\top$ are the singular values
of $A$.  Thus the relevant population eigenvalues are those of
\[
    G(\alpha,\omega) := A^\top A .
\]
With deterministic strengths $\rho_i>0$ and coverage weights
$w_i(\omega)=K\pi_i(\omega)\rho_i^2$, the Gram matrix has the arrowhead form
\[
    G =
    w_0 e_0e_0^\top
    +
    \sum_{i=1}^{r-1}
    w_i
    \left(\sqrt{1-h_i^2}e_0+h_i e_i\right)
    \left(\sqrt{1-h_i^2}e_0+h_i e_i\right)^\top .
\]
Equivalently,
\[
    a=w_0+\sum_{i=1}^{r-1}w_i(1-h_i^2),
    \qquad
    c_i=w_i h_i\sqrt{1-h_i^2},
    \qquad
    d_i=w_i h_i^2.
\]
The eigenvalues are roots of the arrowhead secular equation
\[
    a-\lambda
    -
    \sum_{i=1}^{r-1}\frac{c_i^2}{d_i-\lambda}
    =0,
\]
with the usual interlacing behavior around the diagonal entries $d_i$.
In the implementation, the strengths use
\[
    \rho_i = \eta^{r-1-i},
\]
with default $\eta=1.18$ and experiment-specific overrides only when stated in
the configuration files.

\paragraph{Effective-rank threshold.}
The rank panels report an effective-rank diagnostic,
\[
    \sigma_j > \tau\sigma_1,
    \qquad \tau=0.05,
\]
or equivalently
\[
    \lambda_j(G) > \tau^2\lambda_1(G).
\]
Therefore effective-rank transitions occur at parameter values satisfying
\[
    \lambda_j(G(\alpha,\omega))
    =
    \tau^2\lambda_1(G(\alpha,\omega)).
\]

\paragraph{Heterogeneity and overlap.}
In the heterogeneity sweep, $\omega=0$, so target coverage is fixed and
$\alpha$ changes only the angular spread of the dictionary atoms.  In the
overlap sweep, $\alpha=1$, so the dictionary is already maximally heterogeneous
and $\omega$ changes only the coverage distribution $\pi_i(\omega)$.  This
separates the positive mechanism, angular complementarity, from the redundancy
failure mode, repeated coverage of the same predictive directions.
